% Multiplicity of infection --- plan \S4.7 (Phase 2).
% Sources: new_notes3 Part C \S14 (authoritative, corrected);
% new_notes Part II \S17; new_notes2 \S10.1.
% All formulas numerically verified 2026-08-07 (see file comments in the
% verification notes); the K_2 values below were cross-checked against the
% master equation in new_notes Part II \S17.

\section{Multiplicity of infection}
\label{sec:moi}

So far the cell has been founded by a single particle. In applications this
is not always the right picture: a macrophage may be invaded by several
bacteria at once, and in the laboratory the multiplicity of infection is a
control parameter. Nothing essential changes --- the process is still linear
--- but the bookkeeping does, and one plausible guess turns out to be wrong.
Throughout this section, $k$ founders share a single cell and hence a single
catastrophe clock: rupture occurs at rate $\delta$ times the \emph{total}
load, and dumps the whole cell.

\subsection{Fixation functions for $k$ founders}

By the branching property, the $k$ founder lineages evolve independently
until the shared clock rings. No burst occurs by $t$ precisely when no
lineage has burst; internal extinction occurs precisely when every lineage
has cleared internally. Hence
\begin{equation}
I_k(t)=I(t)^k,\qquad
D_k(t)=D(t)^k,\qquad
\Ihat_k(t)=I(t)^k-D(t)^k.
\label{eq:Ihatk}
\end{equation}

\begin{remark}[A guess that fails]
\label{rem:ihatk}
The tempting guess is $\Ihat_k=\Ihat^{\,k}$: each lineage stays productive
with probability $\Ihat$, independently. But ``productive'' for the whole
cell means \emph{not fixed}, and the intersection of the independent events
``lineage $i$ avoids bursting and does not go extinct'' is not the event
``the cell is not fixed'': the cell also stops being productive when every
lineage dies out. Correctly, $\Ihat_k=I^k-D^k$, whereas
$\Ihat^k=(I-D)^k$; already for $k=2$ the two differ by
$I^2-D^2-(I-D)^2=2D(I-D)\neq0$ whenever $\mu>0$. Numerically, at
$(\beta,\mu,\delta)=(1,0.2,0.05)$ and $t=1$: the correct
$\Ihat_2=I^2-D^2\approx0.8458$, while $\Ihat^2\approx0.6542$. When $\mu=0$,
$D\equiv0$ and the two coincide --- again the invisible case.
\end{remark}

\subsection{Moments and the release kernel}
\label{sec:moi-moments}

The moments fall out of the $k$-founder \pgf{} of
\S\ref{sec:pgf-k-founders}, $G_k=G^k$, where $G$ is the single-founder
(defective) \pgf. Differentiating at $z=1$, with $G(1,t)=I(t)$,
$\partial_zG|_{z=1}=J(t)$ and $\partial_z^2G|_{z=1}=K(t)-J(t)$,
\begin{equation}
J_k(t)=k\,I(t)^{k-1}J(t),
\label{eq:Jk}
\end{equation}
and, from the second factorial moment,
\begin{align}
K_k(t)&=\partial_z^2G_k\big|_{z=1}+\partial_zG_k\big|_{z=1}\nonumber\\
&=k(k-1)I(t)^{k-2}J(t)^2+kI(t)^{k-1}\lrb{K(t)-J(t)}+kI(t)^{k-1}J(t)\nonumber\\
&=k\,K(t)I(t)^{k-1}+k(k-1)J(t)^2I(t)^{k-2}.
\label{eq:Kk}
\end{align}
The expected release flux --- the population quantity that will build the
renewal model of \S\ref{sec:renewal} --- is therefore
\begin{equation}
g_k(t)=\delta K_k(t)
=\delta\lrb{k\,K(t)I(t)^{k-1}+k(k-1)J(t)^2I(t)^{k-2}}.
\label{eq:gk}
\end{equation}

\begin{remark}[Correction of the free-sum form]
\label{rem:gk}
An earlier note gave $g_k=\delta\lrb{kK+k(k-1)J^2}$, the second moment of a
free sum of $k$ independent copies. That omits the powers of $I$: the
second-moment contribution of a configuration is counted only while the
shared clock has \emph{not} rung. At $(\beta,\mu,\delta)=(1,0.2,0.05)$,
$t=1$, $k=2$: the free-sum form gives $K_2\approx23.39$; the corrected
\eqref{eq:Kk} gives $K_2\approx22.264$, which agrees with a direct
master-equation computation. The difference is the superlinearity of $g_k$
in $k$ wearing a disguise: even the free-sum form is superlinear, but the
true kernel is superlinear \emph{and} tempered by survival. That a cell's
expected release flux depends superlinearly on its founding multiplicity is
a genuine prediction of the model --- one that classical constant-$p$ models
cannot express, and that low-multiplicity single-cell assays can test.
\end{remark}

\subsection{Lifetime yield from $k$ founders}

The lifetime mean release from $k$ founders is
$V_\infty^{(k)}=\int_0^\infty g_k(t)\,\dd t$. Carrying out the integration
(deferred to Appendix~\ref{app:technical}),
\begin{equation}
V_\infty^{(1)}=V_\infty=\frac{a(1-b)}{a-1},
\end{equation}
and, for $k\ge2$,
\begin{equation}
V_\infty^{(k)}
=(1-b^k)
+\frac{2\beta}{\delta}\Biggl[\frac{1}{k+1}-b^k+\frac{k\,b^{k+1}}{k+1}\Biggr]
+k(k-1)\frac{\beta}{\delta}\int_1^b (x-a)(x-b)\,x^{k-2}\,\dd x.
\label{eq:Vk}
\end{equation}
In the no-death case the answer simplifies dramatically:
\begin{equation}
V_\infty^{(k)}=k+\frac{\beta}{\delta}\qquad(\mu=0).
\label{eq:Vkmu0}
\end{equation}
Indeed, when $\mu=0$ the event type is independent of the load
(\S\ref{sec:chained}), so the number of births before the single catastrophe
is $\mathrm{Geom}_0(s)$ with mean $\rho/s=\beta/\delta$ whatever the
founding load, and the total release is simply $k$ founders plus that many
births. For example, with $\beta=1$, $\delta=0.1$:
$V_\infty^{(k)}=11,12,13,14,15,\ldots$

Two biological points are worth pausing on. First, the yield is \emph{not}
$kV_\infty$: the first catastrophe dumps the entire cell once, so additional
founders add their own expected descendants but not another cell's worth of
release. At $(\beta,\mu,\delta)=(1,0.2,0.05)$: $V_\infty\approx13.99$, so a
naive scaling would predict $2V_\infty\approx27.97$ for two founders; the
true value is $V_\infty^{(2)}\approx17.43$. Second, while the
unconditional yield grows sub-extensively in $k$, the \emph{conditional}
mean burst size given eventual burst, $V_\infty^{(k)}/(1-b^k)$,
\emph{increases} with $k$ ($17.23,\,18.07,\,19.02,\,21.00$ for
$k=1,2,3,5$ at the same parameters): multiply infected cells are more likely
to burst at all, and when they do, they burst bigger.
